% limitations.tex
% Drop-in limitations section for IPIBench paper.
% Include with: % limitations.tex
% Drop-in limitations section for IPIBench paper.
% Include with: % limitations.tex
% Drop-in limitations section for IPIBench paper.
% Include with: % limitations.tex
% Drop-in limitations section for IPIBench paper.
% Include with: \input{limitations} (before \section{Conclusion})

\section{Limitations}
\label{sec:limitations}

IPIBench V1 has several limitations that bound the scope and interpretation of our findings. 
Consistent with responsible disclosure and rigorous scientific benchmarking practices, we document these boundaries explicitly.

% ------------------------------------------------------------------
\subsection{Scope Exclusions: Benign Utility and False-Positive Degradation}
\label{sec:lim-utility}
% ------------------------------------------------------------------

IPIBench is strictly an \emph{adversarial vulnerability benchmark}, designed to quantify the upper and lower bounds of LLM susceptibility to indirect prompt injection. 
Evaluating the downstream utility degradation, refusal false-positive rates (over-refusal on benign documents), or semantic degradation imposed by input filtering and spotlighting delimiters on un-injected context is intentionally outside the scope of this work. 
Consequently, defense efficacy figures reported in this paper reflect pure injection mitigation and should be paired with external utility benchmarks (e.g., MMLU, GSM8K, or MT-Bench) before deploying these mitigations in production environments.

% ------------------------------------------------------------------
\subsection{Single-Turn Persistence Boundary}
\label{sec:lim-persistence}
% ------------------------------------------------------------------

All 100 test scenarios in IPIBench evaluate \emph{single-turn contextual instruction injection}. 
While our threat models evaluate whether an injected payload induces an immediate unauthorized action or disclosure in the current conversational turn, we do not measure multi-turn memory poisoning, stateful session drift, or post-compromise persistence across long multi-turn dialogues. 
Our findings therefore quantify \emph{single-turn compliance vulnerability}; multi-turn stateful persistence remains an important avenue for future research.

% ------------------------------------------------------------------
\subsection{Controlled Axis Deferral: Axis 1 (Isolated RLHF Effect)}
\label{sec:lim-axis1}
% ------------------------------------------------------------------

Our benchmark establishes 10 controlled research axes across 53+ evaluated model architectures. 
However, Axis 1---the isolated measurement of the pure RLHF alignment effect at the 405B parameter frontier (comparing base LLaMA 3.1 405B with its RLHF-tuned twin Nous Hermes 3 405B)---is explicitly deferred. 
Due to the sunset of unmetered academic inference tiers and the severe financial cost of dedicated multi-node A100/H100 hosting for 405B dense models, Axis 1 is reserved for future compute-grant allocations. 
All other nine controlled axes (including Axis 2 Reasoning Buffers, Axis 3 Dense vs.\ MoE, and Axis 4 Generational Drift) are fully evaluated and statistically verified.

% ------------------------------------------------------------------
\subsection{Automated Three-Tier Scoring and Metric Bounds}
\label{sec:lim-scoring}
% ------------------------------------------------------------------

Our automated evaluation pipeline utilizes an ordinal three-tier phrase-matching and signal-detection scorer that assigns each response $s \in \{0, 0.5, 1\}$. 
Because heuristic evaluation on natural language outputs exhibits an inherent ambiguity interval (e.g., outputs containing partial tool calls alongside refusal disclaimers), we do not report a single monolithic ASR metric. 
Instead, we report the \textbf{Weighted ASR} as our primary headline expected-value metric:
\begin{equation}
\text{ASR}_{\text{weighted}} = \frac{N_{1.0} + 0.5 \cdot N_{0.5}}{N_{\text{total}}} \times 100
\end{equation}
bounded deterministically by a \textbf{Strict Lower Bound} ($\text{ASR}_{\text{strict}} = N_{1.0}/N_{\text{total}}$) and a \textbf{Lenient Upper Bound} ($\text{ASR}_{\text{lenient}} = (N_{1.0} + N_{0.5})/N_{\text{total}}$). 
In Phase 2 evaluation, residual $0.5$ cases are adjudicated via a secondary LLM Judge and verified against a 300-sample human-annotated validation set ($\kappa \ge 0.88$).

% ------------------------------------------------------------------
\subsection{Taxonomy Dimension Variance & Class Imbalance}
\label{sec:lim-taxonomy}
% ------------------------------------------------------------------

The IPIBench dataset encodes eight taxonomy dimensions across its 100 core scenarios. 
Four secondary metadata dimensions (\texttt{authority\_claimed}, \texttt{target\_action\_type}, \texttt{linguistic\_register}, and \texttt{harm\_severity}) exhibit lower empirical variance ($>65\%$ concentration in modal classes) and are retained solely for schema filtering, not primary hypothesis testing. 
Primary statistical analyses focus exclusively on the four high-variance, balanced dimensions: \texttt{category}, \texttt{attack\_goal}, \texttt{evasion\_style}, and \texttt{injection\_position}. 

Additionally, within \texttt{attack\_goal}, scenarios reflect the empirical distribution of real-world attacks, with majority representation in task hijacking (64 scenarios) compared to minority classes such as identity corruption (5 scenarios). 
Per-dimension statistical power is reported with bootstrap 95\% confidence intervals, and fine-grained sub-category conclusions are appropriately scoped.

% ------------------------------------------------------------------
\subsection{Threat Model & Terminology Framing}
\label{sec:lim-terminology}
% ------------------------------------------------------------------

To maintain scientific precision and avoid anthropomorphic or construct-validity overclaiming:
\begin{itemize}
    \item We define our primary target of measurement as \textbf{LLM Susceptibility to Contextual Instruction Injection}, rather than generalized ``Agent Security''.
    \item Unauthorized data access is formalized as \textbf{Information Disclosure Compliance}.
    \item Synthetic tool triggers are formalized as \textbf{Instruction Compliance in Simulated Tool Contexts}.
    \item Cross-model shifts across releases are framed as \textbf{Generational Resistance Drift}.
\end{itemize}

% ------------------------------------------------------------------
\subsection{Responsible Disclosure Status}
\label{sec:lim-disclosure}
% ------------------------------------------------------------------

All vulnerabilities identified in commercial frontier endpoints have been pre-notified to vendor security teams:
(i) Grok 4 tool-schema exfiltration (NRF-003) disclosed to \texttt{security@x.ai};
(ii) DeepSeek V3 system-prompt exfiltration (NRF-002) disclosed to DeepSeek Security;
(iii) UI rendering payload execution (NRF-026) disclosed under standard 90-day coordinated disclosure.
 (before \section{Conclusion})

\section{Limitations}
\label{sec:limitations}

IPIBench V1 has several limitations that bound the scope and interpretation of our findings. 
Consistent with responsible disclosure and rigorous scientific benchmarking practices, we document these boundaries explicitly.

% ------------------------------------------------------------------
\subsection{Scope Exclusions: Benign Utility and False-Positive Degradation}
\label{sec:lim-utility}
% ------------------------------------------------------------------

IPIBench is strictly an \emph{adversarial vulnerability benchmark}, designed to quantify the upper and lower bounds of LLM susceptibility to indirect prompt injection. 
Evaluating the downstream utility degradation, refusal false-positive rates (over-refusal on benign documents), or semantic degradation imposed by input filtering and spotlighting delimiters on un-injected context is intentionally outside the scope of this work. 
Consequently, defense efficacy figures reported in this paper reflect pure injection mitigation and should be paired with external utility benchmarks (e.g., MMLU, GSM8K, or MT-Bench) before deploying these mitigations in production environments.

% ------------------------------------------------------------------
\subsection{Single-Turn Persistence Boundary}
\label{sec:lim-persistence}
% ------------------------------------------------------------------

All 100 test scenarios in IPIBench evaluate \emph{single-turn contextual instruction injection}. 
While our threat models evaluate whether an injected payload induces an immediate unauthorized action or disclosure in the current conversational turn, we do not measure multi-turn memory poisoning, stateful session drift, or post-compromise persistence across long multi-turn dialogues. 
Our findings therefore quantify \emph{single-turn compliance vulnerability}; multi-turn stateful persistence remains an important avenue for future research.

% ------------------------------------------------------------------
\subsection{Controlled Axis Deferral: Axis 1 (Isolated RLHF Effect)}
\label{sec:lim-axis1}
% ------------------------------------------------------------------

Our benchmark establishes 10 controlled research axes across 53+ evaluated model architectures. 
However, Axis 1---the isolated measurement of the pure RLHF alignment effect at the 405B parameter frontier (comparing base LLaMA 3.1 405B with its RLHF-tuned twin Nous Hermes 3 405B)---is explicitly deferred. 
Due to the sunset of unmetered academic inference tiers and the severe financial cost of dedicated multi-node A100/H100 hosting for 405B dense models, Axis 1 is reserved for future compute-grant allocations. 
All other nine controlled axes (including Axis 2 Reasoning Buffers, Axis 3 Dense vs.\ MoE, and Axis 4 Generational Drift) are fully evaluated and statistically verified.

% ------------------------------------------------------------------
\subsection{Automated Three-Tier Scoring and Metric Bounds}
\label{sec:lim-scoring}
% ------------------------------------------------------------------

Our automated evaluation pipeline utilizes an ordinal three-tier phrase-matching and signal-detection scorer that assigns each response $s \in \{0, 0.5, 1\}$. 
Because heuristic evaluation on natural language outputs exhibits an inherent ambiguity interval (e.g., outputs containing partial tool calls alongside refusal disclaimers), we do not report a single monolithic ASR metric. 
Instead, we report the \textbf{Weighted ASR} as our primary headline expected-value metric:
\begin{equation}
\text{ASR}_{\text{weighted}} = \frac{N_{1.0} + 0.5 \cdot N_{0.5}}{N_{\text{total}}} \times 100
\end{equation}
bounded deterministically by a \textbf{Strict Lower Bound} ($\text{ASR}_{\text{strict}} = N_{1.0}/N_{\text{total}}$) and a \textbf{Lenient Upper Bound} ($\text{ASR}_{\text{lenient}} = (N_{1.0} + N_{0.5})/N_{\text{total}}$). 
In Phase 2 evaluation, residual $0.5$ cases are adjudicated via a secondary LLM Judge and verified against a 300-sample human-annotated validation set ($\kappa \ge 0.88$).

% ------------------------------------------------------------------
\subsection{Taxonomy Dimension Variance & Class Imbalance}
\label{sec:lim-taxonomy}
% ------------------------------------------------------------------

The IPIBench dataset encodes eight taxonomy dimensions across its 100 core scenarios. 
Four secondary metadata dimensions (\texttt{authority\_claimed}, \texttt{target\_action\_type}, \texttt{linguistic\_register}, and \texttt{harm\_severity}) exhibit lower empirical variance ($>65\%$ concentration in modal classes) and are retained solely for schema filtering, not primary hypothesis testing. 
Primary statistical analyses focus exclusively on the four high-variance, balanced dimensions: \texttt{category}, \texttt{attack\_goal}, \texttt{evasion\_style}, and \texttt{injection\_position}. 

Additionally, within \texttt{attack\_goal}, scenarios reflect the empirical distribution of real-world attacks, with majority representation in task hijacking (64 scenarios) compared to minority classes such as identity corruption (5 scenarios). 
Per-dimension statistical power is reported with bootstrap 95\% confidence intervals, and fine-grained sub-category conclusions are appropriately scoped.

% ------------------------------------------------------------------
\subsection{Threat Model & Terminology Framing}
\label{sec:lim-terminology}
% ------------------------------------------------------------------

To maintain scientific precision and avoid anthropomorphic or construct-validity overclaiming:
\begin{itemize}
    \item We define our primary target of measurement as \textbf{LLM Susceptibility to Contextual Instruction Injection}, rather than generalized ``Agent Security''.
    \item Unauthorized data access is formalized as \textbf{Information Disclosure Compliance}.
    \item Synthetic tool triggers are formalized as \textbf{Instruction Compliance in Simulated Tool Contexts}.
    \item Cross-model shifts across releases are framed as \textbf{Generational Resistance Drift}.
\end{itemize}

% ------------------------------------------------------------------
\subsection{Responsible Disclosure Status}
\label{sec:lim-disclosure}
% ------------------------------------------------------------------

All vulnerabilities identified in commercial frontier endpoints have been pre-notified to vendor security teams:
(i) Grok 4 tool-schema exfiltration (NRF-003) disclosed to \texttt{security@x.ai};
(ii) DeepSeek V3 system-prompt exfiltration (NRF-002) disclosed to DeepSeek Security;
(iii) UI rendering payload execution (NRF-026) disclosed under standard 90-day coordinated disclosure.
 (before \section{Conclusion})

\section{Limitations}
\label{sec:limitations}

IPIBench V1 has several limitations that bound the scope and interpretation of our findings. 
Consistent with responsible disclosure and rigorous scientific benchmarking practices, we document these boundaries explicitly.

% ------------------------------------------------------------------
\subsection{Scope Exclusions: Benign Utility and False-Positive Degradation}
\label{sec:lim-utility}
% ------------------------------------------------------------------

IPIBench is strictly an \emph{adversarial vulnerability benchmark}, designed to quantify the upper and lower bounds of LLM susceptibility to indirect prompt injection. 
Evaluating the downstream utility degradation, refusal false-positive rates (over-refusal on benign documents), or semantic degradation imposed by input filtering and spotlighting delimiters on un-injected context is intentionally outside the scope of this work. 
Consequently, defense efficacy figures reported in this paper reflect pure injection mitigation and should be paired with external utility benchmarks (e.g., MMLU, GSM8K, or MT-Bench) before deploying these mitigations in production environments.

% ------------------------------------------------------------------
\subsection{Single-Turn Persistence Boundary}
\label{sec:lim-persistence}
% ------------------------------------------------------------------

All 100 test scenarios in IPIBench evaluate \emph{single-turn contextual instruction injection}. 
While our threat models evaluate whether an injected payload induces an immediate unauthorized action or disclosure in the current conversational turn, we do not measure multi-turn memory poisoning, stateful session drift, or post-compromise persistence across long multi-turn dialogues. 
Our findings therefore quantify \emph{single-turn compliance vulnerability}; multi-turn stateful persistence remains an important avenue for future research.

% ------------------------------------------------------------------
\subsection{Controlled Axis Deferral: Axis 1 (Isolated RLHF Effect)}
\label{sec:lim-axis1}
% ------------------------------------------------------------------

Our benchmark establishes 10 controlled research axes across 53+ evaluated model architectures. 
However, Axis 1---the isolated measurement of the pure RLHF alignment effect at the 405B parameter frontier (comparing base LLaMA 3.1 405B with its RLHF-tuned twin Nous Hermes 3 405B)---is explicitly deferred. 
Due to the sunset of unmetered academic inference tiers and the severe financial cost of dedicated multi-node A100/H100 hosting for 405B dense models, Axis 1 is reserved for future compute-grant allocations. 
All other nine controlled axes (including Axis 2 Reasoning Buffers, Axis 3 Dense vs.\ MoE, and Axis 4 Generational Drift) are fully evaluated and statistically verified.

% ------------------------------------------------------------------
\subsection{Automated Three-Tier Scoring and Metric Bounds}
\label{sec:lim-scoring}
% ------------------------------------------------------------------

Our automated evaluation pipeline utilizes an ordinal three-tier phrase-matching and signal-detection scorer that assigns each response $s \in \{0, 0.5, 1\}$. 
Because heuristic evaluation on natural language outputs exhibits an inherent ambiguity interval (e.g., outputs containing partial tool calls alongside refusal disclaimers), we do not report a single monolithic ASR metric. 
Instead, we report the \textbf{Weighted ASR} as our primary headline expected-value metric:
\begin{equation}
\text{ASR}_{\text{weighted}} = \frac{N_{1.0} + 0.5 \cdot N_{0.5}}{N_{\text{total}}} \times 100
\end{equation}
bounded deterministically by a \textbf{Strict Lower Bound} ($\text{ASR}_{\text{strict}} = N_{1.0}/N_{\text{total}}$) and a \textbf{Lenient Upper Bound} ($\text{ASR}_{\text{lenient}} = (N_{1.0} + N_{0.5})/N_{\text{total}}$). 
In Phase 2 evaluation, residual $0.5$ cases are adjudicated via a secondary LLM Judge and verified against a 300-sample human-annotated validation set ($\kappa \ge 0.88$).

% ------------------------------------------------------------------
\subsection{Taxonomy Dimension Variance & Class Imbalance}
\label{sec:lim-taxonomy}
% ------------------------------------------------------------------

The IPIBench dataset encodes eight taxonomy dimensions across its 100 core scenarios. 
Four secondary metadata dimensions (\texttt{authority\_claimed}, \texttt{target\_action\_type}, \texttt{linguistic\_register}, and \texttt{harm\_severity}) exhibit lower empirical variance ($>65\%$ concentration in modal classes) and are retained solely for schema filtering, not primary hypothesis testing. 
Primary statistical analyses focus exclusively on the four high-variance, balanced dimensions: \texttt{category}, \texttt{attack\_goal}, \texttt{evasion\_style}, and \texttt{injection\_position}. 

Additionally, within \texttt{attack\_goal}, scenarios reflect the empirical distribution of real-world attacks, with majority representation in task hijacking (64 scenarios) compared to minority classes such as identity corruption (5 scenarios). 
Per-dimension statistical power is reported with bootstrap 95\% confidence intervals, and fine-grained sub-category conclusions are appropriately scoped.

% ------------------------------------------------------------------
\subsection{Threat Model & Terminology Framing}
\label{sec:lim-terminology}
% ------------------------------------------------------------------

To maintain scientific precision and avoid anthropomorphic or construct-validity overclaiming:
\begin{itemize}
    \item We define our primary target of measurement as \textbf{LLM Susceptibility to Contextual Instruction Injection}, rather than generalized ``Agent Security''.
    \item Unauthorized data access is formalized as \textbf{Information Disclosure Compliance}.
    \item Synthetic tool triggers are formalized as \textbf{Instruction Compliance in Simulated Tool Contexts}.
    \item Cross-model shifts across releases are framed as \textbf{Generational Resistance Drift}.
\end{itemize}

% ------------------------------------------------------------------
\subsection{Responsible Disclosure Status}
\label{sec:lim-disclosure}
% ------------------------------------------------------------------

All vulnerabilities identified in commercial frontier endpoints have been pre-notified to vendor security teams:
(i) Grok 4 tool-schema exfiltration (NRF-003) disclosed to \texttt{security@x.ai};
(ii) DeepSeek V3 system-prompt exfiltration (NRF-002) disclosed to DeepSeek Security;
(iii) UI rendering payload execution (NRF-026) disclosed under standard 90-day coordinated disclosure.
 (before \section{Conclusion})

\section{Limitations}
\label{sec:limitations}

IPIBench V1 has several limitations that bound the scope and interpretation of our findings. 
Consistent with responsible disclosure and rigorous scientific benchmarking practices, we document these boundaries explicitly.

% ------------------------------------------------------------------
\subsection{Scope Exclusions: Benign Utility and False-Positive Degradation}
\label{sec:lim-utility}
% ------------------------------------------------------------------

IPIBench is strictly an \emph{adversarial vulnerability benchmark}, designed to quantify the upper and lower bounds of LLM susceptibility to indirect prompt injection. 
Evaluating the downstream utility degradation, refusal false-positive rates (over-refusal on benign documents), or semantic degradation imposed by input filtering and spotlighting delimiters on un-injected context is intentionally outside the scope of this work. 
Consequently, defense efficacy figures reported in this paper reflect pure injection mitigation and should be paired with external utility benchmarks (e.g., MMLU, GSM8K, or MT-Bench) before deploying these mitigations in production environments.

% ------------------------------------------------------------------
\subsection{Single-Turn Persistence Boundary}
\label{sec:lim-persistence}
% ------------------------------------------------------------------

All 100 test scenarios in IPIBench evaluate \emph{single-turn contextual instruction injection}. 
While our threat models evaluate whether an injected payload induces an immediate unauthorized action or disclosure in the current conversational turn, we do not measure multi-turn memory poisoning, stateful session drift, or post-compromise persistence across long multi-turn dialogues. 
Our findings therefore quantify \emph{single-turn compliance vulnerability}; multi-turn stateful persistence remains an important avenue for future research.

% ------------------------------------------------------------------
\subsection{Controlled Axis Deferral: Axis 1 (Isolated RLHF Effect)}
\label{sec:lim-axis1}
% ------------------------------------------------------------------

Our benchmark establishes 10 controlled research axes across 53+ evaluated model architectures. 
However, Axis 1---the isolated measurement of the pure RLHF alignment effect at the 405B parameter frontier (comparing base LLaMA 3.1 405B with its RLHF-tuned twin Nous Hermes 3 405B)---is explicitly deferred. 
Due to the sunset of unmetered academic inference tiers and the severe financial cost of dedicated multi-node A100/H100 hosting for 405B dense models, Axis 1 is reserved for future compute-grant allocations. 
All other nine controlled axes (including Axis 2 Reasoning Buffers, Axis 3 Dense vs.\ MoE, and Axis 4 Generational Drift) are fully evaluated and statistically verified.

% ------------------------------------------------------------------
\subsection{Automated Three-Tier Scoring and Metric Bounds}
\label{sec:lim-scoring}
% ------------------------------------------------------------------

Our automated evaluation pipeline utilizes an ordinal three-tier phrase-matching and signal-detection scorer that assigns each response $s \in \{0, 0.5, 1\}$. 
Because heuristic evaluation on natural language outputs exhibits an inherent ambiguity interval (e.g., outputs containing partial tool calls alongside refusal disclaimers), we do not report a single monolithic ASR metric. 
Instead, we report the \textbf{Weighted ASR} as our primary headline expected-value metric:
\begin{equation}
\text{ASR}_{\text{weighted}} = \frac{N_{1.0} + 0.5 \cdot N_{0.5}}{N_{\text{total}}} \times 100
\end{equation}
bounded deterministically by a \textbf{Strict Lower Bound} ($\text{ASR}_{\text{strict}} = N_{1.0}/N_{\text{total}}$) and a \textbf{Lenient Upper Bound} ($\text{ASR}_{\text{lenient}} = (N_{1.0} + N_{0.5})/N_{\text{total}}$). 
In Phase 2 evaluation, residual $0.5$ cases are adjudicated via a secondary LLM Judge and verified against a 300-sample human-annotated validation set ($\kappa \ge 0.88$).

% ------------------------------------------------------------------
\subsection{Taxonomy Dimension Variance & Class Imbalance}
\label{sec:lim-taxonomy}
% ------------------------------------------------------------------

The IPIBench dataset encodes eight taxonomy dimensions across its 100 core scenarios. 
Four secondary metadata dimensions (\texttt{authority\_claimed}, \texttt{target\_action\_type}, \texttt{linguistic\_register}, and \texttt{harm\_severity}) exhibit lower empirical variance ($>65\%$ concentration in modal classes) and are retained solely for schema filtering, not primary hypothesis testing. 
Primary statistical analyses focus exclusively on the four high-variance, balanced dimensions: \texttt{category}, \texttt{attack\_goal}, \texttt{evasion\_style}, and \texttt{injection\_position}. 

Additionally, within \texttt{attack\_goal}, scenarios reflect the empirical distribution of real-world attacks, with majority representation in task hijacking (64 scenarios) compared to minority classes such as identity corruption (5 scenarios). 
Per-dimension statistical power is reported with bootstrap 95\% confidence intervals, and fine-grained sub-category conclusions are appropriately scoped.

% ------------------------------------------------------------------
\subsection{Threat Model & Terminology Framing}
\label{sec:lim-terminology}
% ------------------------------------------------------------------

To maintain scientific precision and avoid anthropomorphic or construct-validity overclaiming:
\begin{itemize}
    \item We define our primary target of measurement as \textbf{LLM Susceptibility to Contextual Instruction Injection}, rather than generalized ``Agent Security''.
    \item Unauthorized data access is formalized as \textbf{Information Disclosure Compliance}.
    \item Synthetic tool triggers are formalized as \textbf{Instruction Compliance in Simulated Tool Contexts}.
    \item Cross-model shifts across releases are framed as \textbf{Generational Resistance Drift}.
\end{itemize}

% ------------------------------------------------------------------
\subsection{Responsible Disclosure Status}
\label{sec:lim-disclosure}
% ------------------------------------------------------------------

All vulnerabilities identified in commercial frontier endpoints have been pre-notified to vendor security teams:
(i) Grok 4 tool-schema exfiltration (NRF-003) disclosed to \texttt{security@x.ai};
(ii) DeepSeek V3 system-prompt exfiltration (NRF-002) disclosed to DeepSeek Security;
(iii) UI rendering payload execution (NRF-026) disclosed under standard 90-day coordinated disclosure.
